% !TEX program = lualatex
% La documentation du paquet titrages. Compiler par lualatex (deux passes).
\documentclass[a4paper,11pt]{article}
\usepackage{fontspec}
\usepackage[french]{babel}
\usepackage[margin=2.3cm]{geometry}
\usepackage{amsmath}
\usepackage{siunitx}
\usepackage[version=4]{mhchem}
\sisetup{locale=FR}
\usepackage{booktabs,array}
\usepackage{titrages}
\pgfplotsset{compat=1.18}
\usepackage[most]{tcolorbox}
\tcbuselibrary{listings}
\lstset{basicstyle=\ttfamily\small, columns=fullflexible, keepspaces, breaklines,
  literate={é}{{\'e}}1 {è}{{\`e}}1 {à}{{\`a}}1 {ê}{{\^e}}1}
\newtcblisting{exemple}{colback=gray!4, colframe=gray!50, listing above text,
  left=2mm, right=2mm}
\newcommand{\cs}[1]{\texttt{\textbackslash#1}}
\newcommand{\cle}[1]{\texttt{#1}}
\usepackage[hidelinks]{hyperref}

\title{\textsf{titrages} — courbes de titrage sans approximation\\pH-métrie, conductimétrie, potentiométrie, dans \textsf{pgfplots}}
\author{antnardo\thanks{\url{https://github.com/antnardo/titrages}}}
\date{version 0.1.0 — 23 septembre 2026}

\begin{document}
\maketitle

\begin{abstract}
Le paquet trace les courbes pH $=f(V)$, $\sigma=f(V)$ et $E=f(V)$ du titrage d'un bécher quelconque par une burette quelconque — acides et bases forts ou faibles, polyacides, ampholytes, mélanges, titrant faible compris, couples rédox à plusieurs étapes ou dimériques — avec la dilution, sans aucune approximation de réaction prépondérante.
Sous Lua\LaTeX, il donne aussi le pH, la conductivité ou le potentiel à un volume donné, les points d'inflexion, les équivalences stœchiométriques et la construction de la méthode des tangentes.
Sous pdf\LaTeX, un second moteur trace les courbes avec \textsf{pgfmath} seul.
Il accompagne le paquet Python \textsf{titrages}, qui ajoute une base de couples et d'indicateurs.
\end{abstract}

\tableofcontents

\section{Le principe}

Chercher le pH d'une composition donnée, c'est chercher la racine d'un polynôme de degré $N+2$, sans formule dès que $N>2$.
Le sens inverse est au contraire explicite : \emph{à pH imposé, tout se calcule directement}.
Pour une solution qui contient des espèces faibles $i$, décrites comme des polyacides $\mathrm{H}_{N_i}\mathrm{A}_i$ introduits sous leur forme $q_i$ (le nombre de protons déjà cédés), à la concentration $C_i$, et un excès $b$ de base forte sur l'acide fort, l'électroneutralité s'écrit $R(\mathrm{pH})=0$ avec
\[
R(\mathrm{pH}) = \sum_i C_i\,\bigl(Y_{1,i}(h) - q_i\bigr) - b + \delta(h),
\qquad \delta = \frac{K_e}{h} - h,
\qquad Y_1 = \sum_{j} j\,x_j .
\]
Les fractions $x_j$ sont explicites en $h$, et $R$ est strictement croissante.
Verser un volume $V$ de burette dans $V_0$ de bécher mélange les deux solutions ; l'électroneutralité du mélange, multipliée par $V_0+V$, se sépare en
\[
V_0\,R_\text{bécher}(\mathrm{pH}) + V\,R_\text{burette}(\mathrm{pH}) = 0
\quad\Longrightarrow\quad
V(\mathrm{pH}) = -V_0\,\frac{R_\text{bécher}(\mathrm{pH})}{R_\text{burette}(\mathrm{pH})} .
\]
Le volume est une fonction explicite du pH : la courbe se trace comme une courbe paramétrée par le pH, entre le pH initial (numérateur nul) et celui de la solution titrante (dénominateur nul, asymptote).
Pour un monoacide titré par une base forte, on retrouve l'abscisse normalisée $n = (Y_1 + \delta/C_A)/(1 - \delta/C_B)$.
Le dénominateur $1-\delta/C_B$ est la correction de dilution : tracer $n_0 = Y_1 + \delta/C_A$ comme une courbe de titrage, ainsi que le fait Kalka (2021), c'est compter la base versée sans l'eau qui l'accompagne, et surestimer le pH de $\log 3\approx\num{0,48}$ à deux fois l'équivalence pour des solutions de même concentration.
Les tests du paquet Python confrontent cette formule à un calcul indépendant, qui dilue les concentrations à la main et résout l'électroneutralité du mélange.

Les fractions sont calculées en échelle logarithmique : la forme naïve $x_j = x_0\,k_j/h^j$ fait intervenir $10^{\pm 56}$ pour l'EDTA à pH 14, ce qui dépasse la précision de \textsf{pgfmath} et approche les limites de la double précision.

Les points d'inflexion, ce qu'on mesure, sont calculés à la précision de la machine par le moteur \textsf{lua} : ce sont les minima de $|\mathrm{d}V/\mathrm{d\,pH}|$ (ou $|\mathrm{d}V/\mathrm{d}E|$) sur la courbe directe --- en $V$, le saut de $\mathrm{Fe^{2+}}$ par $\mathrm{MnO_4^-}$ tient en $10^{-9}$\,mL ---, affinés par le zéro de la dérivée de leur logarithme plutôt que par une section dorée, qui plafonne à $\sqrt\varepsilon$. La variance des formes s'écrit $\sum_{j<k}x_jx_k(k-j)^2$, et la pente en dérivant $V_0R_\text{bécher}+VR_\text{burette}=0$ : aucune soustraction de deux nombres voisins au cœur du saut. Les tests les confrontent à un calcul en précision arbitraire : $10^{-10}$ unité de pH, $10^{-11}$\,V.

\section{Installation et chargement}

Les deux fichiers \texttt{titrages.sty} et \texttt{titrages.lua} vont ensemble, dans un dossier que \TeX{} trouve (\texttt{TEXMFHOME/tex/latex/titrages/}, ou le dossier du document).
\begin{lstlisting}
\usepackage{titrages}                % lua sous LuaLaTeX, pgf ailleurs
\usepackage[moteur=pgf]{titrages}    % forcer pgfmath, même sous LuaLaTeX
\end{lstlisting}
Le paquet charge \textsf{pgfplots}. Les tracés s'emploient dans un environnement \texttt{axis}.

\section{Décrire un titrage}

\begin{lstlisting}
\titragedefinir{<nom>}{<clés>}
\end{lstlisting}
enregistre un titrage sous un nom, que les commandes de tracé et de calcul reprennent.

\subsection{Les clés du titrage}

\begin{tabular}{@{}lll@{}}
\toprule
Clé & Défaut & Rôle \\
\midrule
\cle{volume} & 10 & le volume $V_0$ du bécher \\
\cle{becher} & vide & la solution du bécher (clés ci-dessous) \\
\cle{burette} & vide & la solution titrante \\
\cle{volume max} & voir ci-contre & le volume versé maximal \\
\cle{points} & 300 & le nombre de points de la courbe \\
\cle{pKe} & 14 & le produit ionique de l'eau \\
\cle{dilution} & \cle{true} & \cle{false} : la burette apporte ses espèces, pas son eau \\
\cle{abscisse} & \cle{volume} & \cle{n} : les volumes divisés par la première équivalence (lua) \\
\cle{echelle} & 1 & diviser les volumes par ce nombre \\
\bottomrule
\end{tabular}

\medskip
Les volumes s'écrivent dans l'unité qu'on veut, la même partout.
Sans \cle{volume max}, le moteur lua trace jusqu'à la dernière équivalence augmentée de la première ; le moteur pgf, jusqu'à $2V_0$.

\subsection{Les clés d'une solution}

\begin{tabular}{@{}ll@{}}
\toprule
Clé & Rôle \\
\midrule
\cle{espece = \{pKa = <liste>, C = <mol/L>, forme = <q>\}} & une espèce faible, répétable \\
\cle{acide fort = <mol/L>} & un acide fort, cumulable \\
\cle{base forte = <mol/L>} & une base forte, cumulable \\
\bottomrule
\end{tabular}

\medskip
Une espèce est décrite par \emph{tous} les pKa de son couple, dans n'importe quel ordre, et par la forme sous laquelle on l'introduit : \cle{forme = 0} (défaut) est la forme la plus acide, \cle{forme = N} la plus basique.
L'ammoniac est \cle{\{pKa = 9.25, forme = 1\}}, le carbonate de sodium \cle{\{pKa = \{6.35, 10.33\}, forme = 2\}}, l'hydrogénocarbonate \cle{forme = 1}.
Les ions spectateurs n'ont pas à être décrits : une solution préparée est neutre, et leur charge se déduit de la forme.
Les nombres passent par \texttt{\textbackslash fp\_eval:n} : \cle{C = 1e-3} et \cle{C = 10\^{}(-2.5)} sont acceptés.

\section{Tracer}

{\small\begin{tabular}{@{}>{\raggedright\arraybackslash}p{7.8cm}>{\raggedright\arraybackslash}p{5.4cm}l@{}}
\toprule
Commande & Trace & Moteur \\
\midrule
\cs{addplottitrage[<options>]\{<nom>\}} & la courbe pH $=f(V)$ & lua, pgf \\
\cs{addplottitragederivee[<options>]\{<nom>\}} & la pente $\mathrm{d\,pH}/\mathrm{d}V$, exacte & lua \\
\cs{addplottitragesauts[<options>]\{<nom>\}} & les points d'inflexion & lua \\
\cs{addplottitragetangentes}\allowbreak\texttt{[<options>]\{<nom>\}\{<V>\}} & la méthode des tangentes depuis $V$ & lua \\
\cs{addplotfractions[<options>]\{<pKa>\}} & les fractions de chaque forme selon le pH & lua, pgf \\
\bottomrule
\end{tabular}}

\medskip
Les options sont celles de \cs{addplot}.
Les courbes suivent le cycle de styles de \textsf{pgfplots}, sans marques.

\begin{exemple}
\titragedefinir{ethanoique}{
  volume = 10,
  becher = { espece = { pKa = 4.76, C = 0.1 } },
  burette = { base forte = 0.1 },
  volume max = 20,
}
\begin{tikzpicture}
\begin{axis}[width=9cm, height=6cm, xlabel={$V$ (mL)}, ylabel=pH,
             ymin=0, ymax=14, grid=both]
  \addplottitrage[thick]{ethanoique}
  \addplottitragetangentes[red, densely dashed]{ethanoique}{7}
  \addplottitragesauts[black, mark=*]{ethanoique}
\end{axis}
\end{tikzpicture}
\end{exemple}

\subsection{Polyacides, et une équivalence qui ne se voit pas}

L'acide phosphorique a trois acidités, donc trois équivalences stœchiométriques, mais la troisième ($\mathrm{pK_a} = 12{,}35$) ne donne aucun saut : la courbe n'a que deux points d'inflexion.
\begin{exemple}
\titragedefinir{phosphorique}{
  becher = { espece = { pKa = {2.15, 7.20, 12.35}, C = 0.1 } },
  burette = { base forte = 0.1 },
}
\begin{tikzpicture}
\begin{axis}[width=9cm, height=6cm, xlabel={$V$ (mL)}, ylabel=pH, grid=both]
  \addplottitrage[thick]{phosphorique}
  \addplottitragesauts[black, mark=*]{phosphorique}
\end{axis}
\end{tikzpicture}

Équivalences : \titrageequivalence{phosphorique}{1}, \titrageequivalence{phosphorique}{2}
et \titrageequivalence{phosphorique}{3}~mL ; sauts à \titragesaut{phosphorique}{1}
et \titragesaut{phosphorique}{2}~mL.
\end{exemple}

\subsection{Une base faible par un acide fort, en abscisse normalisée}

\begin{exemple}
\titragedefinir{ammoniac}{
  becher = { espece = { pKa = 9.25, C = 0.1, forme = 1 } },
  burette = { acide fort = 0.1 },
  abscisse = n,
}
\begin{tikzpicture}
\begin{axis}[width=9cm, height=5cm, xlabel={$n = V/V_\mathrm{eq}$}, ylabel=pH]
  \addplottitrage[thick]{ammoniac}
\end{axis}
\end{tikzpicture}
\end{exemple}

\subsection{Un mélange, par un titrant faible}

Un bécher d'acide chlorhydrique et d'acide éthanoïque, titré par l'ammoniac : rien ne change dans la description.
\begin{exemple}
\titragedefinir{melange}{
  volume = 20,
  becher = { acide fort = 0.05, espece = { pKa = 4.76, C = 0.05 } },
  burette = { espece = { pKa = 9.25, C = 0.1, forme = 1 } },
  volume max = 30,
}
\begin{tikzpicture}
\begin{axis}[width=9cm, height=5cm, xlabel={$V$ (mL)}, ylabel=pH, grid=both]
  \addplottitrage[thick]{melange}
  \addplottitragesauts[black, mark=*]{melange}
\end{axis}
\end{tikzpicture}
\end{exemple}

\subsection{La dérivée}

\begin{exemple}
\begin{tikzpicture}
\begin{axis}[width=9cm, height=5cm, xlabel={$V$ (mL)},
             ylabel={$\mathrm{d\,pH}/\mathrm{d}V$}, ymode=log]
  \addplottitragederivee[thick]{ethanoique}
\end{axis}
\end{tikzpicture}
\end{exemple}

\subsection{Les fractions}

\begin{exemple}
\titragesreglages{pH min=0, pH max=14}
\begin{tikzpicture}
\begin{axis}[width=9cm, height=4.5cm, xlabel=pH, ylabel=fraction,
             legend pos=outer north east]
  \addplotfractions[thick]{6.35, 10.33}
  \legend{\ce{H2CO3}, \ce{HCO3-}, \ce{CO3^2-}}
\end{axis}
\end{tikzpicture}
\end{exemple}

\section{Conductimétrie}

La loi de Kohlrausch, en solution idéale, donne $\sigma=\sum_i\lambda_i^\circ c_i$.
À pH imposé, toutes les concentrations sont explicites, et le volume aussi : la courbe $\sigma(V)$ est donc, elle aussi, une courbe paramétrée par le pH.
La conductivité demande ce que le pH ignore : la conductivité molaire de chaque forme chargée, et les \emph{ions spectateurs}.

{\small\begin{tabular}{@{}>{\raggedright\arraybackslash}p{5.2cm}>{\raggedright\arraybackslash}p{10.2cm}@{}}
\toprule
Clé & Rôle \\
\midrule
\cle{lambda = \{<l0>, <l1>, …\}} & dans une espèce : $\lambda^\circ$ de chaque forme, en \unit{mS.m^2.mol^{-1}} (0 ou rien pour une forme neutre) \\
\cle{charge = <z0>} & dans une espèce : la charge de sa forme la plus acide (0 par défaut ; 1 pour \ce{NH4+}) \\
\cle{contre ion = <l>} & dans une espèce : $\lambda^\circ$ du contre-ion monovalent qui l'accompagne (\ce{Na+} de l'éthanoate de sodium) \\
\cle{lambda acide fort}, \cle{lambda base forte} & dans une solution : ceux de l'anion de l'acide fort (7,63, \ce{Cl-}) et du cation de la base forte (5,01, \ce{Na+}) \\
\cle{conductivite corrigee} & dans le titrage : multiplier par $(V_0+V)/V_0$, la correction de dilution de TP \\
\cle{unite} & dans le titrage : \cle{mS/cm} (défaut), \cle{S/m} ou \cle{mS/m} \\
\bottomrule
\end{tabular}}

\medskip
\ce{H3O+} (34,98) et \ce{HO-} (19,86) sont connus du paquet.
Une forme chargée sans $\lambda$ est une erreur, pas un zéro silencieux.

\begin{exemple}
\titragedefinir{conducto}{
  volume = 100,
  becher = { acide fort = 0.01, espece = { pKa = 4.76, C = 0.01, lambda = {0, 4.09} } },
  burette = { base forte = 0.1 },
  volume max = 30,
  conductivite corrigee = true,
}
\begin{tikzpicture}
\begin{axis}[width=9cm, height=5.5cm, xlabel={$V$ (mL)},
             ylabel={$\sigma\,(V_0+V)/V_0$ (mS/cm)}, grid=both]
  \addplottitrageconductivite[thick]{conducto}
\end{axis}
\end{tikzpicture}

À \qty{10}{mL} : \qty{\titrageconductivite{conducto}{10}}{mS/cm}.
\end{exemple}

\section{Oxydoréduction}

C'est la transposition exacte : les électrons remplacent les protons, $E/s$ le pH, avec $s=RT\ln 10/F$.
Un système a ses formes de la plus réduite (0) à la plus oxydée ; chaque étape a son $E^\circ$, son nombre d'électrons $n$ et le nombre de protons $m$ que consomme la réduction.
À pH fixé — le milieu est supposé tamponné —, $E^{\circ\prime}=E^\circ - s\,\frac mn\,\mathrm{pH}$, et les fractions s'écrivent comme celles d'un polyacide.
Le bilan d'électrons remplace l'électroneutralité : $R(E)=\sum_i T_i\,(Y_i(E)-e_{q_i})$, croissant, et $V(E) = -V_0R_\text{bécher}/R_\text{burette}$.

Un couple \emph{dimérique} (\ce{I2}/\ce{I-}, \ce{Cr2O7^2-}/\ce{Cr^3+}, \ce{S4O6^2-}/\ce{S2O3^2-}) n'a pas de $Y(E)$ : sa loi de Nernst porte un carré, et une simple dilution change son potentiel.
Le moteur lua résout alors $E(V)$ par dichotomie sur le bilan du mélange ; le moteur pgf ne les accepte pas.

\begin{lstlisting}
\titrageredoxdefinir{<nom>}{<clés>}
\end{lstlisting}

{\small\begin{tabular}{@{}>{\raggedright\arraybackslash}p{5.6cm}>{\raggedright\arraybackslash}p{9.8cm}@{}}
\toprule
Clé & Rôle \\
\midrule
\cle{volume}, \cle{becher}, \cle{burette}, \cle{volume max}, \cle{points}, \cle{abscisse}, \cle{echelle} & comme pour un titrage acide-base \\
\cle{pH} & le pH du milieu (0 par défaut) \\
\cle{reference} & l'électrode de mesure : \cle{ENH} (défaut), \cle{ECS}, \cle{Ag/AgCl}, ou son potentiel en V \\
\cle{pente nernst} & $s$, 0,05916 V à \qty{25}{\degreeCelsius} \\
\cle{espece = \{E0 = …, n = …, H = …, stoechiometrie = \{a, b\}, C = …, forme = …\}} & une espèce : les $E^\circ$ (V/ENH), $n$ et $m$ de chaque étape, la stœchiométrie $a\,\mathrm{Ox}+n\,e^-=b\,\mathrm{Red}$ d'un couple dimérique, la concentration de l'espèce, et sa forme : \cle{red} (défaut), \cle{ox} ou un numéro \\
\bottomrule
\end{tabular}}

\medskip
\cs{addplottitrage}, \cs{addplottitragesauts}, \cs{addplottitragederivee}, \cs{titrageequivalence}, \cs{titragevolume} et \cs{titragesaut[E]} s'appliquent ; \cs{titrageE\{<nom>\}\{<V>\}} donne le potentiel, par rapport à la référence.
Les réglages \cle{E min} et \cle{E max} (−0,5 et 2 V) bornent la courbe du moteur pgf.

\begin{exemple}
\titrageredoxdefinir{vanadium}{
  volume = 10, pH = 0, reference = ECS,
  becher = { espece = { E0 = {-0.26, 0.34, 1.00}, n = {1, 1, 1}, H = {0, 2, 2}, C = 0.1 } },
  burette = { espece = { E0 = 1.51, n = 5, H = 8, C = 0.02, forme = ox } },
  volume max = 40,
}
\titrageredoxdefinir{iodometrie}{
  volume = 10, reference = ECS,
  becher = { espece = { E0 = 0.54, n = 2, stoechiometrie = {1, 2}, C = 0.05, forme = ox } },
  burette = { espece = { E0 = 0.08, n = 2, stoechiometrie = {1, 2}, C = 0.1 } },
  volume max = 20,
}
\begin{tikzpicture}
\begin{axis}[width=9cm, height=6cm, xlabel={$V$ (mL)}, ylabel={$E$ (V/ECS)}, grid=both]
  \addplottitrage[thick]{vanadium}
  \addplottitragesauts[black, mark=*]{vanadium}
  \addplottitrage[thick]{iodometrie}
\end{axis}
\end{tikzpicture}

Équivalences : \titrageequivalence{vanadium}{1}, \titrageequivalence{vanadium}{2},
\titrageequivalence{vanadium}{3}~mL ; iodométrie : \titrageequivalence{iodometrie}{1}~mL,
à \qty{\titragesaut[E]{iodometrie}{1}}{V}.
\end{exemple}

\section{Les nombres}

Sous le moteur lua, ces commandes sont \emph{développables} : elles s'emploient dans \cs{num}, dans une coordonnée, dans un \cs{pgfmathsetmacro}.
Le nombre de décimales, 2 par défaut, se règle par la clé \cle{chiffres} de \cs{titragesreglages}.

{\small\begin{tabular}{@{}>{\raggedright\arraybackslash}p{7.2cm}>{\raggedright\arraybackslash}p{8.1cm}@{}}
\toprule
Commande & Rend \\
\midrule
\cs{titragepH\{<nom>\}\{<V>\}} & le pH après avoir versé $V$ \\
\cs{titragepHinitial\{<nom>\}} & le pH du bécher \\
\cs{titrageE\{<nom>\}\{<V>\}} & le potentiel après avoir versé $V$ (titrage rédox) \\
\cs{titrageconductivite\{<nom>\}\{<V>\}} & la conductivité, dans l'unité du titrage \\
\cs{titragevolume\{<nom>\}\{<pH>\}} & le volume qui amène à ce pH (ou ce potentiel) \\
\cs{titrageequivalence\{<nom>\}\{<k>\}} & la $k$-ième équivalence stœchiométrique (titrant fort) \\
\cs{titragesaut[<champ>]\{<nom>\}\{<k>\}} & le $k$-ième point d'inflexion : \cle{volume}, \cle{pH} (ou \cle{E}) ou \cle{pente} \\
\cs{titragetangentes[<champ>]\{<nom>\}\{<V>\}} & le point de la méthode des tangentes depuis $V$ \\
\bottomrule
\end{tabular}}

\begin{exemple}
pH initial \titragepHinitial{ethanoique} ; demi-équivalence
\titragepH{ethanoique}{5} ; équivalence \titragepH{ethanoique}{10}.

Tangentes depuis \qty{7}{mL} : \qty{\titragetangentes{ethanoique}{7}}{mL},
pH \titragetangentes[pH]{ethanoique}{7}.
\end{exemple}

\section{Les réglages}

\cs{titragesreglages\{<clés>\}} : \cle{chiffres} (2), \cle{pH min} (0) et \cle{pH max} (14), qui bornent le domaine des fractions et, pour le moteur pgf, celui des courbes acide-base ; \cle{E min} et \cle{E max}, celui des courbes rédox du moteur pgf.

\section{Les deux moteurs}

Le moteur \textbf{lua} calcule en double précision, dans \texttt{titrages.lua}, et envoie à \textsf{pgfplots} des coordonnées toutes exactes, réparties à longueur d'arc constante : un pas régulier en volume manquerait le saut, un pas régulier en pH laisserait les zones tampons presque vides.
Il résout le pH à un volume donné par dichotomie — la fonction est monotone, la racine unique — et trouve les sauts comme les maxima de la pente.

Le moteur \textbf{pgf} construit, à la définition du titrage, l'expression de $V(\mathrm{pH})$, et la confie à \textsf{pgfplots} en courbe paramétrée, restreinte à $0 \le V \le V_\text{max}$.
Il n'a besoin que de \textsf{pgfplots} et fonctionne sous pdf\LaTeX, mais ne résout aucune équation : pas de nombres, de sauts ni de tangentes, et pas de couple rédox dimérique.
La conductivité s'y trace comme $\bigl(V(\mathrm{pH}), \sigma(\mathrm{pH})\bigr)$, le potentiel comme $\bigl(V(E), E\bigr)$.
Ses points sont réguliers en pH ; \cle{points} les multiplie au besoin.

\section{Avec Python}

Le paquet Python \textsf{titrages} (\url{https://pypi.org/project/titrages}) calcule la même chose — le même algorithme, confronté à celui-ci par ses tests — et y ajoute une base de couples et d'indicateurs : \texttt{solution(\{"NH3": 0.1\})}, \texttt{Titrage(...)}, \texttt{tracer\_titrage(...)}.
Sa fonction \texttt{vers\_pgfplots} écrit une courbe en coordonnées \textsf{pgfplots}.

\section{Limites}

Les solutions sont idéales : l'activité est la concentration, et les conductivités molaires sont les valeurs limites, qui surestiment $\sigma$ dès quelques \unit{mmol/L}.
La température n'intervient que par \cle{pKe} et \cle{pente nernst} ; les pKa, $E^\circ$ et $\lambda^\circ$ sont ceux qu'on donne.
Les couples rédox qui échangent des protons sont pris à pH fixé ; un métal ou un gaz (\ce{Ag}, \ce{H2}) n'est pas un couple de solution.
Un acide fort est traité exactement, comme une dissociation totale, et non comme un couple de pKa très négatif.

\section{Licence}

MIT. Copyright \textcopyright{} 2026 antnardo.

\end{document}
